% Chapter 2: Literature Review
\section{Introduction}

This chapter provides a comprehensive review of relevant literature in three main areas: educational technology, ontologies and semantic web, and computer science education.

\section{Educational Technology and Learning Systems}

\subsection{Adaptive Learning}

Adaptive learning systems have shown promise in personalizing educational experiences. These systems adjust content, pacing, and presentation based on student performance.

\subsection{Interactive Learning Environments}

Interactive learning environments foster engagement and deeper understanding. Research has demonstrated the pedagogical value of interactive interfaces in knowledge acquisition.

\section{Ontologies and Knowledge Representation}

\subsection{Introduction to Ontologies}

Ontologies provide formal specifications of shared conceptualizations. They enable machines to understand and reason about domain-specific knowledge.

\subsection{Semantic Web Technologies}

The semantic web stack, including RDF, OWL, and SPARQL, provides tools for creating and querying knowledge bases. These technologies are foundational to the proposed framework.

\section{Computer Science Education}

\subsection{CS Curriculum Design}

Effective CS curriculum design requires careful consideration of learning outcomes, content organization, and pedagogical strategies.

\subsection{Current Challenges}

Challenge areas include abstraction level, student motivation, and practical application of theoretical concepts.

\section{Gap in Literature}

While substantial research exists in each area individually, there is limited work integrating ontological approaches with interactive CS education systems. This dissertation addresses that gap.
